\section{Proof of Theorem~\ref{Thm:HQCComplete}}

\begin{proof}
    First, a partitioning will be constructed for $\Sigma_{\CoreExt}^0$.
    Since $\Sigma_{\CoreExt}^0 \subseteq \Sigma_{\CatHQC}$, then the following equation holds by \cref{Lem:FuncDefPart}.
    \begin{equation*}
        \Sigma_{\CoreExt}^0
        =
        \Sigma_{\CoreExt}^0 \cap \Sigma_{\CatHQC}
        =
        \left( \Sigma_{\CoreExt}^0 \cap \Sigma_{\textbf{Prim}} \right)
        \sqcup
        \left( \Sigma_{\CoreExt}^0 \cap \Sigma_{\textbf{Ctrl}} \right)
        \sqcup
        \left( \Sigma_{\CoreExt}^0 \cap \Sigma_{\textbf{Pow}} \right)
    \end{equation*}
    Next, it must be shown that $r_n( - )$ is well-defined for each $n \in \mathbb{N}$.
    This means that $r_n( - )$ assigns each $g \in \Sigma_{\CoreExt}^n$ to a unique word $r_n( g )$ with $r_n( g ) \in \Mor( \Core )$.
    The proof follows by induction on $n$.
    \begin{itemize}
    \item \textbf{Base Case}.
          Let $g \in \Sigma_{\CoreExt}^0$.
          Then there are three cases to consider.
          \begin{enumerate}
          \item Assume that $g \in \Sigma_{\CoreExt}^0 \cap \Sigma_{\textbf{Prim}} = \Sigma_{\textbf{Prim}} \subseteq \Sigma_{\Core}^0$.
                Since $r_0( g ) = g \exparr 1$ by definition, then $r_0( g ) \in \Power( \Sigma_{\Core}^0 ) \subseteq \Sigma_{\Core}^1 \subseteq \Mor( \Core )$.
                Then $r_0( g )$ is well-defined.
          \item Assume that $g \in \Sigma_{\CoreExt}^0 \cap \Sigma_{\textbf{Ctrl}}$ or $g \in \Sigma_{\CoreExt}^0 \cap \Sigma_{\textbf{Pow}}$.
                Then $g \not \in \Sigma_{\CoreExt}^0 \cap \Sigma_{\textbf{Prim}}$.
                Since $\Sigma_{\CoreExt}^0 = \Sigma_{\Core} \cup \Sigma_{\textbf{Prim}}$, then $g \in \Sigma_{\Core} \subseteq \Mor( \Core )$
                Since $r_0( g ) = g$ by definition, then $r_0( g )$ is well-defined.
          \end{enumerate}
          In each case, $r_0( g )$ is well-defined.
          Since $g$ was arbitrary, then $r_0( - )$ is well-defined.
    \item \textbf{Inductive Hypothesis}.
          For some $n \in \mathbb{N}$, $r_n: \Sigma_{\CoreExt}^n \to U( \Core )$ is well-defined.
    \item \textbf{Inductive Step}.
          Assume that for some $n \in \mathbb{N}$, $r_n: \Sigma_{\CoreExt}^n \to U( \Core )$ is well-defined.
          Then there exists a well-defined functor $\Reduce_n = P( r_n )$.
          In particular this means that $\Reduce_n( - )$ terminates when applied to each $U \in \Mor( \CoreExt[n] )$.
          Let $g \in \Sigma_{\CoreExt}^{n+1}$.
          There are three cases to consider.
          \begin{enumerate}
          \item Assume that $g \in \Sigma_{\CoreExt}^n$.
                By assumption, $r_n( - )$ is well-defined, which means that $r( g )$ is unique with $r( g ) \in \Mor( \Core )$.
                Since $r_{n+1}( - )$ assigns $g$ to $r_n( g ) \in \Mor( \Core )$, then $r_{n+1}( g )$ is well-defined.
          \item Assume $g \in \Tower( \CoreExt[n] ) \cap \CoreExt[n+1]( k, k )$ for some $k \in \mathbb{N}$.
                Since $\Reduce_n( - )$ terminates when applied to each $U \in \Mor( \CoreExt[n] )$, then it follows from \cref{Thm:ExpandDef} that $\Expand_{n+1}( - )$ terminates when applied to $g$ with $\Expand_{n+1}( g ) \in \Core( k, k )$.
                Since $r_{n+1}( - )$ assigns $g$ to $\Expand_{n+1}( g )$, then $r_{n+1}( g )$ is well-defined.
          \item Assume that $g \in \Power( \CoreExt[n] ) \cap \CoreExt[n+1]( k, k )$ for some $k \in \mathbb{N}$.
                Then there exists a $U \in \CoreExt[n]( k, k )$ and $\alpha \in \mathbb{R}$ such that $g = U \exparr \alpha$.
                Since $\Reduce_n( - )$ terminates when applied to each $U$, then it follows from \cref{Thm:DropDef} that $\Drop_{n+1}( -, - )$ terminates when applied to $( U, \alpha )$ with $\Drop_{n+1}( U, \alpha ) \in \Core( k, k )$.
                Since $r_{n+1}( - )$ assigns $g$ to $\Drop_{n+1}( U, \alpha )$, then $r_{n+1}( g )$ is well-defined.
          \end{enumerate}
          In each case, $r_{n+1}( g )$ is well-defined.
          Since $g$ was arbitrary, then $r_{n+1}( - )$ is well-defined.
          Then the inductive step holds.
    \end{itemize}
    The by the principle of induction, each $r_n( - )$ is well-defined.
    Since $\bigcup_{n=0}^\infty \Sigma_{\CoreExt}^n = \Sigma_{\CatHQC}$, then this defines a signature morphism $r: \Sigma_{\CatHQC} \to U( \Core )$ such that $r( g ) = r_n ( g )$ for each $n \in \mathbb{N}$ and $g \in \Sigma_{\CoreExt}^n$.
    Define $\Reduce = P( r )$.
    It remains to be shown that if $g \in \Sigma_{\CatHQC}$, then $\Reduce( g ) \approx_E g$.
    The proof follows by induction on $\CoreExt[n] \hookrightarrow \CatHQC$.
    \begin{itemize}
    \item \textbf{Base Case}.
          Let $g \in \Sigma_{\CoreExt}^0$.
          Then there are two cases to consider.
          \begin{enumerate}
          \item Assume that $g \in \Sigma_{\CoreExt}^0 \cap \Sigma_{\textbf{Ctrl}}$.
                Then $r_0( g ) = g \exparr 1$.
                Then $r_0( g ) \approx_E g$ by \HqcPowId.
          \item Assume that $g \in \Sigma_{\CoreExt}^0 \cap \Sigma_{\textbf{Ctrl}}$ or $g \in \Sigma_{\CoreExt}^0 \cap \Sigma_{\textbf{Pow}}$.
                Then $r_0( g ) = g$ by definition.
                Since $( \approx_E )$ is reflexive, then $r_0( g ) \approx_E g$.
          \end{enumerate}
          In each case, $\Reduce( g ) = r_0( g ) \approx_E g$.
          Then $\pi_E( \Reduce( g ) ) = \pi_E( g )$.
          Since $g$ was arbitrary, then $\pi_E \circ \Reduce \circ i = \pi_E \circ i$ where $i: \Sigma_{\CoreExt}^0 \to \CoreExt[0]$.
          Then by the universal property of free prop categories, $\pi_E \circ \Reduce = \pi_E$ when restricted to $\CoreExt[0]$.
          Let $U \in \Mor( \CoreExt[0] )$.
          Then $\pi_E( \Reduce( U ) ) = \pi_E( U )$, which implies that $\Reduce( U ) \approx_E U$.
          Since $U$ was arbitrary, then $\Reduce( U ) \approx_E U$ for all $U \in \Mor( \CoreExt[0] )$.
    \item \textbf{Inductive Hypothesis}.
          For some $n \in \mathbb{N}$, if $U \in \Mor( \CoreExt[n] )$, then $\Reduce( U ) \approx_E U$.
    \item \textbf{Inductive Step}.
          Assume that for some $n \in \mathbb{N}$, if $U \in \Mor( \CoreExt[n] )$, then the derivation $\Reduce( U ) \approx_E U$ holds.
          Let $g \in \Sigma_{\CoreExt}^{n+1}$.
          There are three cases to consider.
          \begin{enumerate}
          \item Assume that $g \in \Sigma_{\CoreExt}^n$.
                Then $\Reduce( g ) \approx_E g$ by the inductive hypothesis.
          \item Assume that $g \in \Tower( \CoreExt[n] )$.
                Then $\Expand( g ) = r_{n+1}( g ) = r( g ) = \Reduce( g )$.
                By the inductive hypothesis, $U \approx_E \Reduce( U ) = \Reduce_n( U )$ for each $U \in \Mor( \CoreExt[n] )$.
                It then follows from \cref{Thm:ExpandDef} that $g \approx_E \Expand( g ) = \Reduce( g )$.
          \item Assume that $g \in \Power( \CoreExt[n] )$.
                Then there exists some $U \in \Mor( \CoreExt[n] )$ and $\alpha \in \mathbb{R}$ such that $g = U \exparr \alpha$.
                Then by definition $\Drop( U, \alpha ) = r_{n+1}( g ) = r( g ) = \Reduce( g )$.
                Then by the inductive hypothesis, $U \approx_E \Reduce( U ) = \Reduce_n( U )$.
                It then follows from \cref{Thm:DropDef} that $g = U \exparr \alpha \approx_E \Drop( U, \alpha ) = \Reduce( g )$.
          \end{enumerate}
          In each case, $\Reduce( g ) \approx_E g$, which implies that $\pi_E( \Reduce_{n+1}( g ) ) = \pi_E( g )$.
          Since $g$ was arbitrary, then $\pi_E \circ \Reduce \circ i = \pi_E \circ i$ for $i: \Sigma_{\CoreExt}^{n+1} \hookrightarrow \CoreExt[n+1]$.
          Then by the universal property of free prop categories, $\pi_E \circ \Reduce = \pi_E$ when restricted to $\CoreExt[n+1]$.
          Let $U \in \CoreExt[n+1]$.
          Then $\pi_E( \Reduce( U ) ) = \pi_E( U )$.
          Then $\Reduce( U ) \approx_E U$.
          Since $U$ was arbitrary, then $\Reduce_{n+1}( U ) \approx_E U$ for all $U \in \CoreExt[n+1]$.
          Then the inductive step holds.
    \end{itemize}
    Then by the principle of induction, $\Reduce( U ) \approx_E U$ for each $n \in \mathbb{N}$ and $U \in \Mor( \CoreExt[n] )$.
    This implies that $\Reduce( g ) \approx_E g$ for each $g \in \Sigma_{\CatHQC}$.
    Let $g \in \Sigma_{\CatHQC}$.
    Since $\Sigma_{\CatHQC} = \bigcup_{n=0}^\infty \Sigma_{\CoreExt}^n$, then there exists an $n \in \mathbb{N}$ such that $g \in \Sigma_{\CoreExt}^n \subseteq \Mor( \CoreExt[n] )$.
    Then $\Reduce( g ) \approx_E g$.
    Since $g$ was arbitrary, then $\Reduce( g ) \approx_E g$ for each $g \in \Sigma_{\CatHQC}$ with $\Reduce: \CatHQC \to \Core$.
    Moreover, $( \Sigma_{\Core}, E )$ is complete with respect to $\sem{-}_H$ by \cref{Thm:CoreComplete}.
    Then $( \Sigma_{\CatHQC}, E )$ is complete with respect to $\sem{-}_H$ by \cref{Thm:Completeness2}.
\end{proof}
